\clearpage
\begin{summarybox}
\textbf{一句话核心：} 内角平分线按邻边长度分对边；逆用时，内部位置与比例缺一不可。

\textbf{速查（非退化三角形，$D$ 在 $BC$ 内部）：}
\[
 \angle BAD=\angle DAC
 \quad\Longleftrightarrow\quad
 \frac{BD}{DC}=\frac{AB}{AC}.
\]
\[
 BD=\frac{BC\cdot AB}{AB+AC},\qquad
 DC=\frac{BC\cdot AC}{AB+AC}.
\]
\[
 BD+DC=BC.
\]
在 $AD$ 为角平分线时，$AB=AC$ 恰好对应 $BD=DC$，也恰好对应 $AD\perp BC$。

\textbf{最短证明记忆：} 同一面积比，从 $A$ 看得 $BD/DC$，从 $D$ 看得 $AB/AC$。

\textbf{最后四查：}
\begin{itemize}[leftmargin=*]
\item 三角形是否非退化？给定三边是否满足严格三角形不等式？
\item 分点是否在 $BC$ 内部？内外角是否混用？
\item 比例是否按 $BD$ 对 $AB$、$DC$ 对 $AC$ 对应？
\item 解出后是否正长度、和为 $BC$，并能回代比例？
\end{itemize}
\end{summarybox}
